\section{Week 10 Refresher: Risk Arithmetic and Assurance Logic}

\subsection*{Setting the Scene}
Governance math is not abstract compliance language; it is decision logic under uncertainty. The field evolved from safety engineering (hazards, layered controls, assurance cases) into AI governance artifacts. The critical insight is traceability: every risk claim must be connected to measurable evidence and triggerable controls.

\subsection*{Notation Introduced in This Chapter}
\begin{itemize}[leftmargin=1.5em]
  \item $R$: baseline risk score.
  \item $R_{\text{res}}$: residual risk after controls.
  \item $\eta_i$: efficacy of control $i$.
  \item $m_t$: monitored metric at time $t$.
  \item $m_0$: baseline metric value.
  \item $\epsilon$: escalation threshold.
\end{itemize}

\subsection*{What You Need To Recall Precisely}
\begin{itemize}[leftmargin=1.5em]
  \item \textbf{Risk decomposition:} severity and likelihood are different axes with different uncertainty.
  \item \textbf{Layered controls:} multiplicative residual-risk models assume forms of independence.
  \item \textbf{Trigger logic:} indicator functions implement threshold decisions.
  \item \textbf{Assurance structure:} claim $\to$ evidence $\to$ control action $\to$ artifact.
  \item \textbf{Correlation caveat:} multiplicative risk-reduction models are optimistic under shared-cause failures.
\end{itemize}

\subsection*{How These Results Build on Each Other}
\begin{enumerate}[leftmargin=1.5em]
  \item Define risk factors and their measurement status.
  \item Model baseline and post-control residual risk.
  \item Apply threshold policy to operational metrics.
  \item Attach each decision to auditable evidence.
\end{enumerate}

\subsection*{Reminder Glossary (Term by Term)}
\begin{itemize}[leftmargin=1.5em]
  \item \textbf{Hazard:} undesirable outcome class to prevent/mitigate.
  \item \textbf{Control efficacy:} estimated fractional risk reduction from a control.
  \item \textbf{Residual risk:} remaining risk after controls.
  \item \textbf{Threshold policy:} decision rule on monitored metrics.
  \item \textbf{Assurance artifact:} evidence item supporting a claim.
\end{itemize}

\subsection*{Worked Mini Example (Single vs Two Controls)}
If baseline risk score is $R=1.0$, and control efficacies are $\eta_1=0.4$, $\eta_2=0.3$:
\[
R_1=R(1-\eta_1)=0.6,\qquad
R_{12}=R(1-\eta_1)(1-\eta_2)=0.42.
\]
Interpretation: layered controls reduce residual risk multiplicatively under the model assumptions.

\subsection*{Worked Example (Threshold and Escalation)}
Suppose drift metric baseline is $m_0=0.12$ and policy threshold is $\epsilon=0.05$.  
Current observed $m_t=0.19$ gives deviation $|m_t-m_0|=0.07>\epsilon$.  
Escalation rule triggers high-review mode.  
Reasoning chain: measured metric -> threshold test -> policy action -> logged evidence.

\subsection*{Intuition Checkpoints}
\begin{itemize}[leftmargin=1.5em]
  \item Multiplicative control models can be optimistic if failures are correlated.
  \item Governance thresholds are policy design decisions, not immutable laws.
  \item Audit readiness is about traceable evidence chains, not only numeric scores.
\end{itemize}

\subsection*{Further Refresh}
\begin{itemize}[leftmargin=1.5em]
  \item Metric reliability and uncertainty: see Week 1 refresher.
  \item Evaluation decomposition and harness logic: see Week 9 refresher.
  \item \textbf{Leveson, \emph{Engineering a Safer World}} --- hazard/control argument structure.
  \item \textbf{NIST AI RMF} --- operational governance controls and documentation patterns.
\end{itemize}
